\subsection{Magnetic hyperbolic Laplacian} \label{ssection:landau} By the Gauss-Bonnet formula, the metric being
hyperbolic, the volume is
\begin{equation}
\int_{\Si} \vol = 4 \pi (\textsl{g} -1),
\end{equation}
where $\textsl{g}\geqslant
2 $ is the genus of $\Si$. Let $L \to \Sigma$ be a Hermitian line bundle with
a Hermitian connection
$\nabla$ whose curvature is $-i B \op{vol}
$, $B > 0$ being a positive constant. Since the Chern class of $L$ is $ (2
\pi)^{-1} B \cdot [ \op{vol}]$, the
degree of $L$ is
\begin{equation}
\label{equation:deg}
\mathrm{deg}(L) = (2 \pi)^{-1} B \int_{\Si} \vol = 2 B (\textsl{g}-1) \in \Z.
\end{equation}
Thus, the pair $(L,
\nabla)$ exists if and only if $2 B (\textsl{g}-1)$ is an integer. When it exists, it
is unique up to tensor product by a flat Hermitian line bundle. The
space of flat Hermitian line bundles up to equivalence is in one-to-one
correspondence trough the holonomy representation with the character space
$\op{Mor} (\pi_1 (\Si), \op{U}(1)) \simeq \op{U} (1)^{2\textsl{g}}$. An~example of a pair $(L, \nabla)$ is provided by the canonical bundle $K := T^{1,0} \Si$,
endowed with its Chern connection, whose curvature is $- i \vol$, so~in this
case $B=1$.

Let $\Delta := \frac{1}{2} \nabla^* \nabla $ be the Laplacian acting on $C^\infty
(\Si, L)$. Since $\Si$ is compact, $\Delta$~has a discrete spectrum and each
eigenvalue has finite multiplicity. We~denote by $\la_0 \leqslant \la_1
\leqslant \cdots $ the eigenvalues of $\Delta$. The following result was
proved in \cite{Iengo-Li-94} (see also \cite{Tejero-06, Landau2}):

\begin{proposition} \label{theo:Landau_Level}
The following holds:
\begin{enumerate}

\item For any pair $(L, \nabla)$ with curvature $F_\nabla=-iB
\op{vol}$ with $B>0$, the first $N := \lfloor B \rfloor $ eigenvalues are
given by
\begin{gather} \label{eq:valeur_propre}
\la_{m } = B (\sfrac{1}{2} + m) - \dfrac{ m (m+1) } {2},\qquad 0
\leqslant m < N,
\end{gather}
and the multiplicity of $\la_m$ is $2 (\textsl{g}-1)(B-\sfrac{1}{2} - m)$ when $m \leqslant N- 2$.

\item If $(L,\nabla) = (K^r,\nabla^{\mathrm{Chern}})$ with $ r \geq 1$, then \eqref{eq:valeur_propre} holds for
$m \leqslant r$ and the remaining part of the spectrum is given by
$\la_{r+n} = \la_r + \frac{1}{2} \mu_n$ where $\{ \mu_n, \; n \geq 0 \}$ is the
spectrum of the Laplace-Beltrami operator of $(\Sigma, g)$.
\end{enumerate}
\end{proposition}

We call the eigenvalues $\la_m$ given by \eqref{eq:valeur_propre} and their eigenspaces the Landau levels.
These quantum levels correspond to the classical energy
levels $S^\lambda \Si$ with $\lambda <B$ on which the magnetic flow is periodic. A~first
clue of this fact is that \eqref{eq:valeur_propre} applied with $m = B$ gives $\frac{1}{2} B^2$ which is exactly the classical
energy $\frac{1}{2} \lambda^2$ at the critical value $\lambda =B$.
Moreover, as was noticed by Comtet \cite{Comtet-86}, \eqref{eq:valeur_propre} can be interpreted as a
Bohr-Sommerfeld condition.

On a different perspective, a common characteristic
of the magnetic flow below the critical energy and the Landau levels is
that they do not depend on the choice of hyperbolic metric $g$. The part of the spectrum above the
critical energy has a different nature (as seen with $L = K^r$ for instance).

\begin{proof}[Sketch of proof] This follows from the Riemann-Roch theorem and the
following two identities
\begin{gather} \label{eq:bochner_bosonic}
\Delta = \Box_L + \tfrac{1}{2} B, \qquad \con{\partial}_L
\con{\partial}_L^* = \Box_{L \otimes K^{-1}} +B -1.
\end{gather}
Here, $L$ is equipped with its holomorphic structure whose Chern connection
is $\nabla$, $\con{\partial}_L: C^{\infty} (L)
\rightarrow C^{\infty} (L \otimes K^{-1})$ is the $d$-bar
operator with the identification $\con{ K} = K^{-1}$ induced by
the metric, and $\Box_L = \con{\partial}_L^*
\con{\partial}_L$. The first identity of \eqref{eq:bochner_bosonic} is a
classical Bochner identity, the second one is proved in \cite[Th.\,7.1]{Landau1}.

For any $m \in \N$, let $\con{\partial}_m = \con{\partial}_{L \otimes
K^{-m}}$ and $\Box_m = \Box_{L \otimes K^{-m}}$. The curvature of $L
\otimes K^{-m}$ being $\frac{1}{i} (B -m) \op{vol}$, the second identity
gives
\[
\con{\partial}_m \Box_m = (
\Box_{m+1} + (B - m) - 1) \con{\partial}_m.\]
Introduce the operator $\Box_{m}^{-1}$
of $C^\infty(L \otimes K^{-m})$ equal to $0$ on $\ker \Box_m$
and inverting $\Box_m$ on the orthogonal of $\ker \Box_m$. Then
$\Box_m^{-1} \con{\partial}_m^*$ is the inverse of $\con{\partial}_m$ on
$(\ker \Box_m)^{\perp}$, so~by the previous equation
\begin{gather*}
\Box_m = \Box_m^{-1} \con{\partial}_m^* \bigl(\Box_{m+1} + B - (m+1)
\bigr) \con{\partial}_m.
\end{gather*}
As a consequence, the spectra satisfy
\begin{gather} \label{eq:rec}
\op{Sp} (\Box_m) \setminus \{ 0 \} = \op{Sp} (
\Box_{m+1}) + B - (m +1).
\end{gather}
Moreover, since the degree of $L \otimes K^{-m}$ is $(B- m) 2 (\textsl{g} -1)$, by
the Riemann-Roch theorem, $h^0 (L \otimes K^{-m}) := \op{dim} (\ker \Box_m)$ is
positive when $B - m \geqslant 1$, and equal to $(B-m) 2 (\textsl{g}-1) + 1 - \textsl{g}$ when
$B- m \geqslant 2 $.

We deduce the first part of the theorem as follows: by the first identity of
\eqref{eq:bochner_bosonic}, $\Delta = \Box_0 + \frac{1}{2} B$ and if $B
\geqslant 1$, $\la_0 = \tfrac{1}{2} B$ and its multiplicity is $h^0 (
L)$. If $B\geqslant 2$, then $H^{0} (\Sigma, L \otimes K^{-1}) >0$, so
$\la_1 =\frac{1}{2}B + B -1$ by \eqref{eq:rec} and its multiplicity is $H^{0}
(\Sigma, L \otimes K^{-1})$. We~can repeat this until we can no longer apply the Riemann-Roch theorem.

When $L = K^r$, after $r$ iterations, it~comes that $\op{Sp}(\Delta) \setminus \{ \la_0, \ldots, \la_{r-1} \}$
is equal to the spectrum of $\Box_r + \frac{1}{2} r^2$. By Bochner identity,
$\Box_r$ is the Laplace-Beltrami operator.
\end{proof}

\section{Twisted semiclassical calculus}

\label{section:microlocal}

We now introduce the semiclassical limit we will be working with. In~this section, $\Sigma$ need just be a closed manifold.
Let $L \to \Sigma$ be a Hermitian line bundle with a Hermitian connection $\nabla$. We~do not
make any assumption on the curvature $F_{\nabla}$. This section is organized as follows:
\begin{itemize}

\item In Section~\ref{ssection:uniform-pseudo}, we~introduce the twisted semiclassical pseudodifferential calculus;

\item In Section~\ref{ssection:defect-measures}, we~discuss semiclassical defect measures associated to quasimodes in this calculus.
\end{itemize}

\subsection{Definition and main properties}

\label{ssection:uniform-pseudo}

For $m \in
\R$, denote by $\Psi^m_{\mathrm{sc}}(\Sigma)$ the space of standard
($h$-dependent) semiclassical pseudodifferential operators of degree $m \in
\R$ (see \cite[Ch.\,14]{Zworski-12} for an introduction). We~emphasize that
$A \in \Psi^m_{\mathrm{sc}}(\Sigma)$ is a \emph{family} of operators $A =
(A_h)_{h > 0}$, where $A_h = \Op_h(a_h) + \mc{O}(h^\infty)$, $\Op_h$ is an
arbitrary semiclassical quantization on $\Sigma$, and $a_h \in
S^m_h(T^*\Sigma)$ satisfies uniform symbolic estimates in $h > 0$ (see \cite[Ch.\,3]{Lefeuvre-book} for a definition of symbol classes). In~addition, the family $A$ might not be defined for all values of $h > 0$, but only on a subset.

\begin{definition}[Twisted semiclassical quantization]
A family of operators $\bA := (\bA_k)_{k \geq 0}$ such that
\[
\bA_{k} : C^\infty(\Sigma, L^{k}) \to C^\infty(\Sigma, L^{k}),
\]
belongs to the space of \emph{twisted} semiclassical operators $\Psi^m_{\mathrm{tsc}}(\Sigma)$ of degree $m \in \R$ if the following holds: for all contractible open subset $U \subset \Sigma$, for all $\chi,\chi' \in C^\infty_{{\comp}}(U)$, for all $s \in C^\infty(U,L)$ such that $|s|=1$ fiberwise, there exists a family of semiclassical operators $A := (A_h) \in \Psi^m_{\mathrm{sc}}(\Sigma)$ such that for all $f \in C^\infty_{{\comp}}(U)$:
\begin{equation}
\label{equation:def}
s^{-k}\bA_k(f s^{k}) = A_{1/k} f.
\end{equation}
\end{definition}

This quantization was introduced by the first author \cite{Charles-00}. We
also refer to \cite[\S 3.2]{Cekic-Lefeuvre-24} and \cite[\S 2]{Charles--23} for a general introduction to this quantization procedure. We~now recall some of its properties.

Define on the open set $U$ the real-valued $1$-form $\beta \in C^\infty(U,T^*U)$ by $\nabla s = -i \beta \otimes s$ (hence $\nabla s^{k} = -i k \beta \otimes s^{k}$). We~also introduce the symplectic $2$-form
\begin{equation}
\label{equation:omega}
\Omega := dA + i \pi^* F_\nabla,
\end{equation}
where $dA$ is the Liouville $2$-form on $T^*\Sigma$ (see \eqref{equation:liouville}) and $\pi : T^*\Sigma \to \Sigma$ is the projection. (For the magnetic Laplacian on surfaces, we~will have $\Omega = dA + B \pi^*\vol$.) It can be easily checked that \eqref{equation:omega} defines a non-degenerate (closed) $2$-form. Given $p \in C^\infty(T^*\Sigma)$, the Hamiltonian vector field of $p$ associated to $\Omega$ is denoted by $H_p^\Omega$. It is defined through the relation
\[
dp = \Omega(\csbullet, H_p^\Omega).
\]
Finally, the Poisson bracket of two observables $p,q \in C^\infty(T^*M)$ computed with respect to $\Omega$ is written as $\{p,q\}^\Omega =: H^\Omega_p q$. We~sum up the main properties of this calculus (below, $S^m_k(T^*\Sigma)$ denotes the space of symbols of order $m$ which may depend on $k \geq 0$; however, the symbolic estimates are required to be uniform in $k$):

\begin{proposition}
The following properties hold:

\begin{enumerate}

\item \emph{Algebra property.} The space
\[\textstyle
\Psi^\csbullet_{\mathrm{tsc}}(\Sigma) := \bigcup_{m \in \R} \Psi^m_{\mathrm{tsc}}(\Sigma)
\]
is a graded algebra. Namely, for all $\bA\in\Psi^m_{\mathrm{tsc}}(\Sigma)$, $\bB\in\Psi^{m'}_{\mathrm{tsc}}(\Sigma)$, $\bA \circ \bB \in \Psi^{\color{black}m+m'}_{\mathrm{tsc}}(\Sigma)$.

\item \emph{Principal symbol.} For $\bA \in \Psi^m_{\mathrm{tsc}}(\Sigma)$, there is a well-defined principal symbol given by
\[
\sigma_{\bA}(x,\xi ; k) = a(x,\xi + \beta(x)) \in
S^m_{k}(T^*\Sigma)/k^{-1}S^{m-1}_{k}(T^*\Sigma),
\]
where $a$ is the principal symbol of $A_{1/k}$.

\item \emph{Commutator.} For $\bA \in \Psi^m_{\mathrm{tsc}}(\Sigma)$, $\bB \in \Psi^{m'}_{\mathrm{tsc}}(\Sigma)$, $[\bA,\bB] \in k^{-1}\Psi^{m+m'-1}_{\mathrm{tsc}}(\Sigma)$ with principal symbol
\[
\sigma_{k[\bA,\bB]} = -i \{\sigma_{\bA},\sigma_{\bB}\}^\Omega.
\]

\item \emph{Calderon-Vaillancourt.} Let $\bA \in \Psi^0_{\mathrm{tsc}}(\Sigma)$. Then $\bA_k : L^2(\Sigma,L^k) \to L^2(\Sigma,L^k)$ is bounded for all $k \geq 0$ with norm
\[
\|\bA_k\|_{L^2 \to L^2} \leq \|\sigma_{\bA}(\csbullet;k)\|_{L^\infty(T^*\Sigma)} + \mc{O}(k^{-1}).
\]

\end{enumerate}

\end{proposition}

By construction, the operator $\Delta_k := \tfrac{1}{2} {\color{black}k^{-2}} (\nabla^k)^*\nabla^k \in \Psi^2_{\mathrm{tsc}}(\Sigma)$ belongs to this calculus and has principal symbol $p(x,\xi) := \sigma_{\Delta_k}(x,\xi) = \tfrac{1}{2}|\xi|^2_g$.

Conversely, one defines a quantization map $\Op : S^m_{k}(T^*\Sigma) \to \Psi^m_{\mathrm{tsc}}(\Sigma)$ as follows. Consider a smooth function $\chi \in C^\infty(\Sigma \times \Sigma)$, equal to $1$ near the diagonal $\Delta \subset \Sigma \times \Sigma$ and supported in $\{d(x,y) < \iota(g)/2026\}$, where $\iota(g)$ is the injectivity radius of the metric $g$. For $k \geq 0$, and $x,y \in \mathrm{supp}(\chi)$, let $\tau_{x \to y} : L^k_x \to L^k_y$ denote the parallel transport with respect to $\nabla$ along the unique $g$-geodesic joining $x$ to $y$. We~then set:
\[
\Op_{k}(a_{k}) f(x) := \dfrac{1}{2\pi} \int_{\Sigma} \int_{T^*_x\Sigma} e^{-ik\xi(\exp^{-1}_x(y))} a_{k}(x,\xi) \tau_{y \to x}(f(y)) \chi(x,y) \dd\xi \dd y,
\]
where $\dd y$ is the Riemannian volume, and $\dd\xi$ the induced volume in the fibers of $T^*\Sigma$. It is a standard calculation to verify that
\[
\sigma_{\Op_{k}(a_{k})} = [a_{k}] \in S^m_{k}(T^*\Sigma)/k^{-1}S^{m-1}_{k}(T^*\Sigma).
\]

An operator $\bA \in \Psi^\csbullet_{\mathrm{tsc}}(\Sigma)$ is \emph{microlocally compactly supported} if $\bA = \Op_k(a_k) + \mc{O}(k^{-\infty})$ where $a_k \in C^\infty_{\comp}(T^*\Sigma)$ has (uniformly in $k$) compact support on $T^*\Sigma$, and the remainder term $\mc{O}(k^{-\infty})$ denotes an operator with smooth Schwartz kernel on $\Sigma \times \Sigma$, all of whose derivatives are $\mc{O}(k^{-\infty})$ in $L^\infty$-norm. We~denote by $\Psi^{\comp}_{\mathrm{tsc}}(\Sigma)$ the set of microlocally compactly supported operators.

Finally, we~shall use the radial compactification $\overline{T^*\Sigma} :=
T^*\Sigma \sqcup \partial_\infty T^*\Sigma$ of cotangent space (obtained by
adding a sphere at infinity to each fiber of $T^*\Sigma$), see \cite[App.\,E.1.3]{Dyatlov-Zworski-book}. As in standard semiclassical theory, an operator
$\bA \in \Psi^m_{\mathrm{tsc}}(\Sigma)$ is \emph{elliptic} at $(x_0,\xi_0)
\in \overline{T^*\Sigma}$ if there exists a constant $c > 0$ such that for all
$k \geq 0$ large enough, $|\sigma_{\bA}(x,\xi;k)| \geqslant
c\langle\xi\rangle^m$ in a neighborhood of $(x_0, \xi_0)$. The \emph{elliptic set} of an operator is denoted by $\Ell(\bA) \subset \overline{T^*\Sigma}$ and the \emph{characteristic set} is the complement of the elliptic set. Operators can be inverted modulo $\mc{O}(k^{-\infty})$ smoothing operators on their elliptic set (parametrix construction).

\subsection{Defect measures}

\label{ssection:defect-measures}

Quasimodes of energy $E \geq 0$ are sections $u_{k} \in C^\infty(\Sigma,L^k)$ satisfying
\begin{equation}
\label{equation:qm}
(k^{-2}\Delta_k-E) u_{k} = \mc{O}_{L^2}(\eps_{k}), \qquad \|u_k\|^2_{L^2}=1,
\end{equation}
where $\eps_{k} \to_{k \to \infty} 0$. Up to extraction along a subsequence $(k_n)_{n \geq 0}$, there exists a measure $\mu$ on $T^*\Sigma$ such that for all $a \in C^\infty_{{\comp}}(T^*\Sigma)$,
\begin{equation}
\label{equation:limit}
\lim_{n \to \infty} \langle\Op_{k_n}(a)u_{k_n},u_{k_n}\rangle_{L^2} = \int_{T^*\Sigma} a(x,\xi) d\mu(x,\xi).
\end{equation}
That $\Delta_k$ is elliptic on $\partial_\infty T^*\Sigma$ (boundary at
infinity of the radial compactification of~$T^*\Sigma$), implies that $\mu$ is
a probability measure, see \cite[App.\,E, Exer.\,23]{Dyatlov-Zworski-book}. In~the following, recall that $p(x,\xi) = \tfrac{1}{2}|\xi|^2_g$.

\begin{proposition}
\label{proposition:scl-measure}
Let $\mu$ be a probability measure associated to the quasimodes \eqref{equation:qm}. Then the following holds:
\begin{enumerate}

\item\label{proposition:scl-measurei} Assume that $\eps_k = o(1)$. Then
\[
\supp(\mu) \subset \mathrm{Char}(k^{-2}\Delta_k-E) = \{p=E\}.
\]

\item\label{proposition:scl-measureii} Assume that $\eps_k = o(1/k)$. Then $\mu$ is invariant by the magnetic Hamiltonian vector field $H_{p}^\Omega$, that is for all $a \in C^\infty_{{\comp}}(T^*\Sigma)$,
\[
\int_{T^*\Sigma} H_{p}^\Omega a(x,\xi) ~\dd\mu(x,\xi) = 0.
\]
\end{enumerate}
\end{proposition}

We refer to \cite[App.\,E.3]{Dyatlov-Zworski-book} for a proof.

\subsection{Technical estimates}

We will need some precise estimates on the remainder in the standard operations of semiclassical analysis (product, commutator, propagation, etc.). Recall that $(\Phi_t)_{t \in \R}$ is the magnetic flow on $T^*\Sigma$, that is the Hamiltonian flow generated by $H_p^\Omega$.

\begin{proposition} For any $E>0$, there exists a constant $C > 0$ such that for all $k \geq 0$, for all $a,b \in
C^\infty_{\comp}(T^*\Sigma)$ with $\mathrm{supp}(a), \mathrm{supp}(b) \subset
\{|\xi|_g \leq E\}$, the following holds:
\begin{align}
\label{equation:calderon-fine}
\|\Op_k(a)\|_{L^2 \to L^2} & \leq \|a\|_{L^\infty(T^*\Sigma)} + k^{-1/2}\|a\|_{C^{14}(T^*\Sigma)}, \\
\label{equation:product-fine}
\|\Op_k(a)\Op_k(b)-\Op_k(ab)\|_{L^2 \to L^2} &\leq Ck^{-1} \|a\|_{C^{17}(T^*\Sigma)}\|b\|_{C^{17}(T^*\Sigma)}, \\
\label{equation:commutator-fine}
\|[\Op_k(a),\Op_k(b)]-\Op_k(-ik^{-1}\{a,b\})\|_{L^2 \to L^2} &\leq Ck^{-1}\|a\|_{C^{17}(T^*\Sigma)}\|b\|_{C^{17}(T^*\Sigma)}, \\
\label{equation:egorov-fine}
\|e^{it k^{-1} \Delta} \Op_k(a) e^{-itk^{-1}\Delta}-\Op_k(a \circ \Phi_t)\|_{L^2 \to L^2} &\leq C k^{-1} \int_0^t \|a \circ \Phi_s\|_{C^{17}(T^*\Sigma)} \dd s.
\end{align}

\end{proposition}

The above estimates are very likely suboptimal by far; however, we~did not try to optimize them.

\begin{proof}
Estimate \eqref{equation:calderon-fine} follows from \cite[Th.\,5.26]{Nonnenmacher-notes} applied in dimension $2$ and Lemma \cite[Lem.\,A.4]{Dyatlov-Jin-Nonnenmacher-22}. The estimates \eqref{equation:product-fine} and \eqref{equation:commutator-fine} can be found in \cite[Lem.\,A.6]{Dyatlov-Jin-Nonnenmacher-22} for the standard semiclassical quantization. It is an exercise to verify that the proofs go through for the twisted quantization. Finally, \eqref{equation:egorov-fine} follows from \eqref{equation:commutator-fine}. Indeed, since $e^{it k^{-1} \Delta}$ is unitary on $L^2(\Sigma, L^k)$, it~suffices to estimate $\|B(t)\|$, where
\[
B(t) := \Op_k(a) - e^{-itk^{-1}\Delta}\Op_k(a\circ\Phi_t)e^{itk^{-1}\Delta}.
\]
Notice that $B(0)=0$ and using \eqref{equation:commutator-fine}:
\[
\begin{split}
B'(t) &= e^{-itk^{-1}\Delta}\left(k[ik^{-2}\Delta_k,\Op_k(a\circ\Phi_t)]-\Op_k(H_p a \circ\Phi_t) \right) e^{itk^{-1}\Delta} \\
& = \mc{O}_{L^2 \to L^2}(\|a\circ\Phi_t\|_{C^{17}(T^*\Sigma)} h),
\end{split}
\]
where $p$ is the principal symbol of $k^{-2}\Delta_k$. Writing
\[
B(t) = \int_0^t B'(s)\dd s,
\]
we find the claimed result.
\end{proof}

\section{Construction of eigenstates}

\label{section:constant}

We now assume that $(\Sigma,g)$ is a Riemannian surface of constant curvature $-1$, genus $\geq 2$, and $L \to \Sigma$ is a Hermitian line bundle equipped with a unitary connection $\nabla$ with curvature $F_\nabla = -iB\vol$ and $B > 0$ is constant. Without loss of generality, since $H^2(\Sigma,\Z) \simeq \Z$ and we will be
considering power $L^k \to \Sigma$, we~can further assume that $\deg(L)=1$.
Note that by \eqref{equation:deg}, this forces $B = |\chi(\Sigma)|^{-1}$. Hence, in what follows, the magnetic intensity on $M$ \emph{is fixed and equal to $B = |\chi(\Sigma)|^{-1}$} on $L \to \Sigma$.

This section is organized as follows:
\begin{itemize}

\item In Section~\ref{ssection:averaging}, we~adapt Weinstein's averaging method \cite{Weinstein-77} to our context and introduce an operator $\bA_k$ whose principal symbol generates a $2\pi$-periodic flow on $\{p < E_c\}$ which is a renormalization of the magnetic flow;

\item We then use it in Section~\ref{ssection:quasimodes} to construct specific eigenstates concentrating in phase space on any closed orbit of the magnetic flow.
\end{itemize}

\subsection{The averaging method} \label{ssection:averaging} By Proposition \ref{theo:Landau_Level}
applied to $L^k \to \Sigma$ with curvature $F_{\nabla^k}=-ikB
\op{vol}$, the first eigenvalues of $\Delta_k = \tfrac{1}{2}\nabla^*\nabla : C^\infty(\Sigma,L^k) \to C^\infty(\Sigma,L^k)$ are given for $m < \lfloor kB\rfloor $ by
\begin{equation}
\label{equation:ev}
\lambda_{k,m} = kB (m+\sfrac{1}{2})-\sfrac{m(m+1)}{2}.
\end{equation}
Notice that $\lambda_{k, \lfloor kB \rfloor - 1 } \sim k^2 E_c$ as $k
\rightarrow \infty$, where $E_c
= \tfrac{1}{2} B^2$ is the critical energy. As a consequence, for any $\eps >0$, we~have that
\[
\op{Sp}
(k^{-2} \Delta_k) \cap [0, E_c -\eps ] \subset \{k^{-2} \la_{k,m} ~|~ m < \lfloor
kB\rfloor \}\]
when $k$ is sufficiently large.

\subsubsection{Weinstein's periodic operator} The following paragraph is
inspired by Weinstein's averaging method \cite{Weinstein-77}. Let $\Pi_{k,m}$ be the orthogonal projector of $C^\infty(M, L^k)$ onto
$\ker(\Delta _k -\lambda_{k,m})$, and set:
\begin{equation}
\label{equation:a}
\bA_k := k^{-1} \sum_{m =0 }^{ \lfloor kB \rfloor -1} m \Pi_{k,m}.
\end{equation}
Observe that by construction $\op{Sp}(\bA_k) \subset k^{-1} \Z_{\geq 0}$ and thus
\begin{equation}
\label{equation:cool}
e^{2i k \pi \bA_k } = \mathbf{1}.
\end{equation}
Furthermore, by~\eqref{equation:ev}, on the space
\[
\mc{I}_k := \bigoplus_{m=0}^{ \lfloor kB \rfloor -1}
\ker(\Delta_k-\lambda_{k,m}),
\]
the operator $\bA_k$ satisfies the identity
\begin{equation}
\label{equation:relation}
k^{-2} \Delta_k = B(\bA_k+ \tfrac{1}{2} k^{-1})- \tfrac{1}{2}
\bA_k(\bA_k+ k^{-1}).
\end{equation}
Equivalently,
\begin{equation}
\label{equation:relation2}
\bA_k= B-\tfrac{1}{2} k^{-1} - \sqrt{ B^2 - 2 k^{-2}
\Delta_k + \tfrac{1}{4} k^{-2}},
\end{equation}
on $\mc{I}_k$. Finally, observe that
\begin{equation}
\label{eq:integrale_projecteur}
\Pi_{k,m} = \dfrac{1}{2\pi} \int_0^{2\pi} e^{-imt} e^{itk\bA_k} \dd t.
\end{equation}
This turns out to be a Fourier Integral Operator in the semiclassical regime $k \to \infty$.

Let
\[
D^*\Sigma := \{(x,\xi) \in T^*\Sigma ~|~ |\xi| < B\}.
\]
Let $a$ be the function defined on $D^*\Sigma$ through the relation
\begin{equation}
\label{equation:mercredi-aprem}
p(x,\xi) = \beta (a (x, \xi)), \qquad \forall (x,\xi) \in D^*\Sigma,
\end{equation}
where $\beta : [0,B] \rightarrow [0, \frac{1}{2} B^2 ]$ is the function $ \beta (s) := B s -\frac{1}{2} s^2$. The following holds:

\begin{lemma}
Let $\chi \in C^\infty_{{\comp}}(T^*\Sigma)$ such that $\supp(\chi) \subset D^*\Sigma$ and $\varphi \in C^\infty(\R)$ such that $\supp(\varphi) \subset (-\infty,B)$. Then
\[
\Op_k(\chi) \bA_k, \varphi(\bA_k) \in \Psi^{\comp}_{\mathrm{tsc}}(\Sigma),
\]
and these operator have respective principal symbols $\chi a$ and $\varphi(a)$.
\end{lemma}

We emphasize that $\bA_k$ is \emph{not} a pseudodifferential operator because of the spectral cutoff involved in its definition (the spectral projector onto $\mc{I}_k$). However, in what follows, we~will always use $\Op_k(\chi) \bA_k, \varphi(\bA_k)$ which \emph{are} pseudodifferential operators. For simplicity, we~also say that $a$ is the principal symbol of $\bA_k$, although this only makes sense in $D^*\Sigma$.

\begin{proof}
We define $f_k : [0,\tfrac{1}{2}B^2] \to \R$ by:
\begin{equation}
\label{equation:fk}
f_k(s) = B-\tfrac{1}{2}k^{-1}-\sqrt{B^2-2s+\tfrac{1}{4}k^{-2}}.
\end{equation}
Notice that $\varphi \circ f_k$ is smooth with compact support in $[0,\tfrac{1}{2}B^2)$ by the assumption made on the support of $\varphi$; it can thus be extended to a smooth compactly supported function $c_k$ on $[0,\infty)$. In~addition, using \eqref{equation:relation2}, the equality $\varphi(\bA_k) =\varphi(f_k(k^{-2}\Delta_k)) = c_k(k^{-2}\Delta_k)$ holds. The function $c_k$ has compact support on $[0,\infty)$ and uniformly bounded derivatives as $k \to \infty$. Hence, by~standard functional calculus (see \cite[Th.\,14.9]{Zworski-12} for instance), $c_k(k^{-2}\Delta_k) = \varphi(\bA_k) \in \Psi^{\comp}_{\mathrm{tsc}}(\Sigma)$.

We now establish that $\Op_k(\chi) \bA_k \in \Psi^{\comp}_{\mathrm{tsc}}(\Sigma)$. For that, we~write
\[
\Op_k(\chi) \bA_k =\Op_k(\chi) \psi(k^{-2}\Delta_k) \bA_k + \Op_k(\chi) (\mathbbm{1}-\psi(k^{-2}\Delta_k)) \bA_k,
\]
where $\psi $ is smooth with compact support in $(-\infty, E_c)$ and such that $\psi(p) \equiv 1$ on the support of $\chi$. Notice that $\Op_k(\chi) \psi(k^{-2}\Delta_k) \bA_k = \Op_k(\chi) d_k(k^{-2}\Delta_k)$, where $d_k(x) = \psi(x) f_k(x)$ and $f_k$ was defined in \eqref{equation:fk}. The function $d_k$ extends to a smooth compactly supported function on $[0,\infty)$ and the same argument as in the previous paragraph using functional calculus shows that $\Op_k(\chi) \psi(k^{-2}\Delta_k) \bA_k \in \Psi^{\comp}_{\mathrm{tsc}}(\Sigma)$. Let us now prove that $\Op_k(\chi) (\mathbbm{1}-\psi(k^{-2}\Delta_k)) \bA_k$ is a residual operator, namely that it has smooth Schwartz kernel all of whose derivatives are $\mc{O}(k^{-\infty})$. By standard arguments, it~suffices to prove that
\[
\Op_k(\chi) (\mathbbm{1}-\psi(k^{-2}\Delta_k)) \bA_k : H^{-N}_k(\Sigma,L^k) \to H^{+N}_k(\Sigma,L^k)
\]
is bounded for all $N \geq 0$ with norm $\mc{O}(k^{-\infty})$, where the space $H^s(\Sigma,L^k)$ is defined as the completion of $C^\infty(\Sigma,L^k)$ with respect to the norm
\begin{equation}
\label{equation:sobolev-norm}
\|u\|^2_{H^s(\Sigma,L^k)} := \|(k^{-2}\Delta_k+\mathbbm{1})^{s/2} u\|^2_{L^2(\Sigma,L^k)}.
\end{equation}
By the support property of $\psi$, the operator $\Op_k(\chi) (\mathbbm{1}-\psi(k^{-2}\Delta_k))$ clearly satisfies this property (this is a composition of two operators with disjoint microsupport). Hence, it~suffice to show that $\bA_k : H^{-N}_k(\Sigma,L^k) \to H^{-N}_k(\Sigma,L^k)$ is bounded with norm $\mc{O}(1)$ {\color{black}for all integers $N \geq 0$}. Using \eqref{equation:sobolev-norm}, this amounts to proving that
\[
(k^{-2}\Delta_k+\mathbbm{1})^{-N/2} \bA_k (k^{-2}\Delta_k+\mathbbm{1})^{N/2} : L^2(\Sigma,L^k) \to L^2(\Sigma,L^k)
\]
has norm $\mc{O}(1)$. However, this is immediate using \eqref{equation:a} and
the fact that this operator acts diagonally on the Fourier decomposition into eigenmodes of $\Delta_k$.

Finally, to compute the principal symbol of these operators, merely observe
that by \eqref{equation:relation}, we~have $p = B a - a^2/2 = \beta(a)$.
\end{proof}

The function $a$ will play an important role in the sequel. One important property is the following:

\begin{lemma}
\label{lemma:2pi-periodic}
The Hamiltonian vector field $H_a^\Omega$ of the symbol $a$ generates a $2 \pi$\nobreakdash-peri\-odic flow on $D^*\Sigma$.
\end{lemma}

\begin{proof}This can be seen by writing $a = \al (p)$ with $\al(y) := B - \sqrt{ B^2 - 2 y }$, and noticing
that $\al' (E) = (B^2 - 2 E)^{-\sfrac{1}{2}}$ is the period of the orbits of $H_p^\Omega$
with energy $p =E$ (see Proposition \ref{proposition:dynamic}, item \eqref{proposition:dynamici}).
\end{proof}

Given a function $f \in C^\infty (D^*\Sigma)$, introduce its average
$\langle f \rangle \in C^\infty (D^*\Sigma)$ by:
\begin{equation}
\label{equation:average}
\langle f \rangle (z) := \frac{1}{2\pi} \int_0^{2 \pi} f (\Phi_t^a (z)) ~\dd t,
\end{equation}
where $(\Phi_t^a)_{t \in \R}$ is the $2\pi$-periodic Hamiltonian flow generated by $H_a^\Omega$ on $D^*\Sigma$:

\begin{lemma} \label{lem:commutateur_espace_propre}
Fix $\eps > 0$. For all $f \in C^\infty_{{\comp}} (D^*\Sigma)$, $k \geq 0$, for all $0 \leq m \leq (1-\eps)kB$, the following holds:
\[
\Pi_{k,m}\Op_{k}(f) \Pi_{k,m} = \Op_{k}(\langle f \rangle)
\Pi_{k,m} + \bigo_{\mc{L}(L^2)}(k^{-1}).
\]
Here the $\bigo $ depends on $f$ but is uniform with respect to $m$.
\end{lemma}

The proof is a variation on Egorov's theorem (see \cite[Th.\,15.2]{Zworski-12} for instance).

\begin{proof}
In the proof, the $\mc{O}$'s are measured in operator norm $L^2 \to L^2$. Fix
$f \in C^\infty_{{\comp}}(D^*\Sigma)$ and choose $\varphi \in
C^\infty(-\infty, B)$
such that $ \varphi (a) =1$ on a neighbourhood of the support of $f$. Then
\[
\Op_k(f) \varphi (\bA_k) \equiv \Op_k(f) \equiv \varphi (\bA_k) \Op_k
(f),
\]
where $\equiv$ stands for equality modulo $\bigo (k^{-\infty})$. This yields
\begin{xalignat*}{2}
[ \Op_k(f), \bA_k ] & \equiv [ \Op_k(f), \varphi (\bA_k) \bA_k ]
= \frac{1}{ik} \Op_k (\{ f, \varphi (a) a \}) + \bigo (k^{-2}) \\ & =
\frac{1}{ik} \Op_k (\{ f, a \}) + \bigo (k^{-2}).
\end{xalignat*}
Write $U_t := e^{itk\bA_k}$. Integrating the previous relation, we~obtain the weak form of Egorov's theorem for $\bA_k$:
\begin{gather} \label{eq:egorov}
U_t \Op_k(f) U_{-t} = \Op_k(f \circ \Phi_t^a) + \bigo (k^{-1}).
\end{gather}
Since $U_t \Pi_{k,m} = \Pi_{k,m} U_t = e^{itm} \Pi_{k,m}$, we~obtain that
\[
\Pi_{k,m} \Op_k(f) \Pi_{k,m} = \Pi_{k,m} U_t \Op_k(f) U_{-t} \Pi_{k,m} = \Pi_{k,m} \Op_k(f \circ
\Phi_t^a) \Pi_{k,m} + \bigo (k^{-1}).
\]
Averaging over time, we~find that
\begin{gather}
\label{eq:moyenne}
\Pi_{k,m} \Op_k(f) \Pi_{k,m} = \Pi_{k,m} \Op_k(\langle f \rangle) \Pi_{k,m} + \bigo (k^{-1}).
\end{gather}
Moreover, since $\langle f \rangle \circ \Phi_t^a = \langle f \rangle $, we
deduce from \eqref{eq:egorov} that $ [ U_t, \Op_k(\langle f \rangle) ] =
\bigo (k^{-1})$ and from \eqref{eq:integrale_projecteur} that $[\Pi_{k,m}, \Op_k(
\langle f \rangle)] = \bigo (k^{-1})$. The claim then follows from this last
relation and \eqref{eq:moyenne}.
\end{proof}
